% =====================================================================
%  ilacion.tex
%  Cadenas analíticas del "Reporte de educción e ilación de requisitos
%  de RedPepper" (secciones 3.3, 4.3, 5.3 y 6.3).
%  Cada cadena muestra el paso evidencia -> interpretación -> requisito.
% =====================================================================

\section{Ilación}
\label{sec:ilacion}

\noindent{\small\textbf{Versión 3} \textperiodcentered{} documento de trabajo sujeto a validación con el interesado.}\par\vspace{0.5em}

La ilación es el razonamiento explícito que relaciona datos de la educción
para identificar necesidades, reglas, conflictos, consecuencias o
requisitos candidatos. Cada cadena indica premisas, operación de análisis,
conclusión y límites. Puede incluir interpretación, contraste y formulación
de hipótesis; una observación aislada no demuestra una regla universal. Su
salida es una conclusión argumentada que debe validarse cuando excede la
evidencia.

Las cuatro tablas siguientes reorganizan las 52 cadenas del reporte de
origen. Cada fila permite revisar el paso \emph{evidencia} (dato de la
educción), \emph{interpretación} (razonamiento) y \emph{requisito}
(candidato derivado o, cuando la cadena no genera obligación, la
conclusión de contexto). Los requisitos se detallan en la
Sección~\ref{sec:especificacion}.


\subsection{Registro de reorganización y justificación}

Las tablas de esta sección reorganizan las cadenas existentes. Mantienen sus
identificadores, evidencias, requisitos candidatos y límites de validación.

\begin{longtable}{L{2.0cm} L{5.6cm} L{6.4cm}}
\caption{Cambios de presentación y razón de cada cambio.}
\label{tab:cambios-ilacion}\\
\toprule
\textbf{Cambio} & \textbf{Aplicación} & \textbf{Razón} \\
\midrule
\endfirsthead
\toprule
\textbf{Cambio} & \textbf{Aplicación} & \textbf{Razón} \\
\midrule
\endhead
\bottomrule
\endfoot
C-01 & Cadenas I-PRE, I-OBS, I-ENT1 e I-ENT2 dispuestas en cuatro tablas. & Permite comparar evidencia, interpretación y resultado sin perder el vínculo con cada evento. \\
C-02 & Requisito y conclusión agrupados en la última columna; límite destacado en la misma celda. & Hace visible qué se propone y qué requiere confirmación antes de aprobarlo. \\
C-03 & La educción separa fuente de negocio, autoría colectiva y participación por verificar. & Evita atribuir a Jerry propuestas del equipo o adjudicar a un integrante una actividad sin respaldo documental. \\
C-04 & Se conserva la tabla de integración entre eventos y sus referencias. & Mantiene la trazabilidad entre datos originales, cadenas y requisitos especificados. \\
\end{longtable}

\subsection{Cadenas del evento EV-PRE-01}
\label{subsec:ilacion-pre}

\begin{longtable}{L{2.35cm} L{2.15cm} L{4.35cm} L{5.15cm}}
\caption{Cadenas de ilación del evento EV-PRE-01 (14 cadenas).}
\label{tab:i-pre}\\
\toprule
\textbf{Cadena} & \textbf{Evidencia} & \textbf{Razonamiento} & \textbf{Resultado y validación} \\
\midrule
\endfirsthead
\toprule
\textbf{Cadena} & \textbf{Evidencia} & \textbf{Razonamiento} & \textbf{Resultado y validación} \\
\midrule
\endhead
\bottomrule
\endfoot
\textbf{I-PRE-01. Identificación de ventas diarias} & E-PRE-01. & la solicitud requiere conservar fecha e importe de cada venta y distinguir registros confirmados de borradores. & \textbf{Requisito:} R-PRE-04. \textbf{Conclusión:} registro de ventas asociado a la jornada. \textbf{Límite:} confirmar si la fecha corresponde al cobro, pedido o entrega. \\
\textbf{I-PRE-02. Costo por producto frente a compra semanal} & E-PRE-02. & se solicita costo directo del producto, mientras el punto 5 usa compras semanales. Son bases distintas y no deben sumarse sin una política que evite duplicación. & \textbf{Requisito:} R-PRE-04 y R-PRE-11. \textbf{Conclusión:} necesidad de registrar costos directos y definir su base de cálculo. \textbf{Límite:} validar si se usan recetas y consumo, costo unitario manual o asignación de compras. \\
\textbf{I-PRE-03. Cierre de jornada} & E-PRE-03. & el total exige un criterio de corte y selección de ventas incluidas. Un resumen no obliga por sí solo a bloquear todas las correcciones. & \textbf{Requisito:} R-PRE-10. \textbf{Conclusión:} consolidado diario con detalle de ventas incluidas. \textbf{Límite:} confirmar anulaciones, devoluciones y horario de corte. \\
\textbf{I-PRE-04. Definición del margen solicitado} & E-PRE-04. & la resta expresa una métrica diaria; sus entradas deben usar la misma jornada y una política de costos definida. & \textbf{Requisito:} R-PRE-04. \textbf{Conclusión:} margen diario con ventas y costos trazables. \textbf{Límite:} la calidad contable depende de cómo se reconoce el costo directo. \\
\textbf{I-PRE-05. Aclaración del objeto de prorrateo} & E-PRE-05. & la frase original dice ventas, aunque se motiva en compras semanales. F-02, sección 5.2, documenta la aclaración posterior de repartir compras. Se conserva la discrepancia y se adopta el reparto de costos como regla candidata apoyada en ese contraste. & \textbf{Requisito:} R-PRE-01, R-PRE-02, R-PRE-03 y R-PRE-04. \textbf{Conclusión:} asignación semanal del costo y cargo diario de adicionales, sin redistribuir ingresos. \textbf{Límite:} obtener confirmación explícita de Jerry sobre el objeto del prorrateo y feriados. \\
\textbf{I-PRE-06. Imputación mensual de gastos} & E-PRE-06. & concepto, monto y mes de imputación permiten consolidar gastos sin confundir fecha de pago y periodo al que corresponden. & \textbf{Requisito:} R-PRE-05. \textbf{Conclusión:} registro de gastos indirectos por periodo. \textbf{Límite:} confirmar devengo, pagos adelantados y gastos compartidos. \\
\textbf{I-PRE-07. Separación de BAI y resultado antes de intereses} & E-PRE-07. & el texto llama BAI a un subtotal anterior a impuestos e intereses. Cuando existen intereses, esas denominaciones no coinciden. La ilación identifica el conflicto; no cambia la declaración de origen. & \textbf{Requisito:} R-PRE-06. \textbf{Conclusión:} subtotales separados y nombres aprobados por el contador. \textbf{Límite:} definir si se desea resultado antes de financiación e impuestos o BAI después de financiación. \\
\textbf{I-PRE-08. Determinación de reglas tributarias} & E-PRE-08. & ventas del mes no basta para definir todos los tributos. Se necesitan régimen, tributo, base, vigencia y aprobación del contador. & \textbf{Requisito:} R-PRE-07. \textbf{Conclusión:} motor de reglas tributarias configurables para el caso. Fuentes de consulta para sus variaciones: N-07 y N-12 a N-14; régimen del negocio no confirmado. \textbf{Límite:} no se infiere una tasa única ni impuesto a la renta siempre proporcional a ventas. \\
\textbf{I-PRE-09. Resultado monetario y porcentaje} & E-PRE-09. & el cliente usa margen y ganancia como equivalentes. Debe distinguirse importe neto del periodo de margen porcentual. & \textbf{Requisito:} R-PRE-07. \textbf{Conclusión:} resultado neto monetario y, si se aprueba, porcentaje sobre ventas. \textbf{Límite:} confirmar denominador y tratamiento de ventas cero. \\
\textbf{I-PRE-10. Marco de presentación} & E-PRE-10. & la intención del cliente requiere una plantilla y clasificaciones aprobadas. No acredita obligación de aplicar NIIF 18 ni cumplimiento integral. & \textbf{Requisito:} R-PRE-08. \textbf{Conclusión:} reporte financiero mensual con plantilla contable identificada. \textbf{Límite:} validar marco aplicable y política de presentación con el contador. \\
\textbf{I-PRE-11. Vínculo entre comprobante y egreso} & E-PRE-11. & guardar archivos sin vínculo no permite explicar qué costo sustentan. Se deriva una asociación recuperable entre documento y registro. & \textbf{Requisito:} R-PRE-09 y R-PRE-14. \textbf{Conclusión:} repositorio documental vinculado a compras y gastos. Contraste documental: N-08 y N-09; una copia adjunta no demuestra por sí sola la deducibilidad fiscal. \textbf{Límite:} acordar formatos, retención, acceso y tratamiento de datos personales. \\
\textbf{I-PRE-12. Catálogo y estadísticas de venta} & E-PRE-12. & el catálogo define qué se vende y a qué precio; el historial evita recalcular ventas antiguas con un precio nuevo. El ranking se basa en cantidades e importes, no prueba margen por plato. & \textbf{Requisito:} R-PRE-12 y R-PRE-13. \textbf{Conclusión:} carta editable, precio aplicado por venta y resúmenes de productos. \textbf{Límite:} confirmar métricas del ranking y autorización para cambios de precio. \\
\textbf{I-PRE-13. Separación de emisión y reporte interno} & E-PRE-13. & el cliente mantiene la emisión de comprobantes de venta en SUNAT. El software solicitado lleva control interno de gastos y su sustento; no debe confundirse con un emisor electrónico. & \textbf{Requisito:} R-PRE-09 y R-PRE-14. \textbf{Conclusión:} reporte diario de egresos vinculado al consolidado mensual. \textbf{Límite:} confirmar eventual importación de comprobantes; no se presupone integración automática con SUNAT. \\
\textbf{I-PRE-14. Conciliación de reporte y soportes} & E-PRE-14. & el esquema exige que los totales del mes puedan desagregarse en registros sustentados. También mantiene la ambigüedad del BAI antes de intereses, identificada en I-PRE-07. & \textbf{Requisito:} R-PRE-06, R-PRE-08 y R-PRE-14. \textbf{Conclusión:} reporte mensual conciliable con gastos diarios y subtotales definidos. \textbf{Límite:} la plantilla gerencial inicial debe conciliarse con el modelo contable aprobado. \\
\end{longtable}

\subsection{Cadenas del evento EV-OBS-01}
\label{subsec:ilacion-obs}

\begin{longtable}{L{2.35cm} L{2.15cm} L{4.35cm} L{5.15cm}}
\caption{Cadenas de ilación del evento EV-OBS-01 (5 cadenas).}
\label{tab:i-obs}\\
\toprule
\textbf{Cadena} & \textbf{Evidencia} & \textbf{Razonamiento} & \textbf{Resultado y validación} \\
\midrule
\endfirsthead
\toprule
\textbf{Cadena} & \textbf{Evidencia} & \textbf{Razonamiento} & \textbf{Resultado y validación} \\
\midrule
\endhead
\bottomrule
\endfoot
\textbf{I-OBS-01. Separación de tareas} & E-OBS-01. & la disposición y los roles descritos permiten distinguir captura del pedido, preparación y entrega. Esta interpretación modela responsabilidades, sin afirmar que cada rol sea una persona distinta. & \textbf{Requisito:} ninguno. Conclusión de contexto: no se especifica una obligación independiente. \textbf{Conclusión:} modelo operativo con responsabilidades de caja, cocina y despacho. \textbf{Límite:} confirmar quién desempeña cada rol por turno. \\
\textbf{I-OBS-02. Pedido como desencadenante} & E-OBS-02. & si la preparación empieza por un pedido entrante, cocina necesita conocer los productos solicitados. La automatización de esa comunicación es una consecuencia propuesta; no fue observada como software ya existente. & \textbf{Requisito:} R-OBS-02. \textbf{Conclusión:} necesidad de entregar a cocina información del pedido confirmado. \textbf{Límite:} la universalidad de la regla debe confirmarse con Jerry; la entrevista 1 aporta corroboración documental. \\
\textbf{I-OBS-03. Separación de demoras} & E-OBS-03. & la espera de cocina y la demora que añada una interfaz tienen causas diferentes. Reducir el tiempo de registro puede evitar espera administrativa, pero no elimina el tiempo de cocción. & \textbf{Requisito:} R-OBS-04. \textbf{Conclusión:} necesidad de captura y envío ágiles del pedido. \textbf{Límite:} no existen mediciones de tiempos ni muestra de pedidos en las fuentes. \\
\textbf{I-OBS-04. Seguimiento de pedidos} & E-OBS-04. & la separación de áreas introduce un punto de transferencia de información. Para evitar confusión al entregar, se propone conservar un identificador común durante el recorrido. & \textbf{Requisito:} R-OBS-01 y R-OBS-03. \textbf{Conclusión:} identificador de pedido y seguimiento de preparación y entrega. \textbf{Límite:} los estados Pendiente, En preparación y Listo son candidatos del equipo. \\
\textbf{I-OBS-05. Límite entre evidencia y solución} & E-OBS-05. & la fuente presenta el POS táctil como solución derivada. No corresponde registrarlo como un dispositivo ya observado ni atribuir un tamaño en píxeles al comportamiento de cocina. & \textbf{Requisito:} ninguno. Conclusión de contexto: no se especifica una obligación independiente. \textbf{Conclusión:} separación entre necesidad de agilidad y alternativa de interfaz. \textbf{Límite:} el umbral mayor a 48 px procede de la especificación previa y debe validarse en el equipo real. \\
\end{longtable}

\subsection{Cadenas del evento EV-ENT-01}
\label{subsec:ilacion-ent1}

\begin{longtable}{L{2.35cm} L{2.15cm} L{4.35cm} L{5.15cm}}
\caption{Cadenas de ilación del evento EV-ENT-01 (5 cadenas).}
\label{tab:i-ent1}\\
\toprule
\textbf{Cadena} & \textbf{Evidencia} & \textbf{Razonamiento} & \textbf{Resultado y validación} \\
\midrule
\endfirsthead
\toprule
\textbf{Cadena} & \textbf{Evidencia} & \textbf{Razonamiento} & \textbf{Resultado y validación} \\
\midrule
\endhead
\bottomrule
\endfoot
\textbf{I-ENT1-01. Necesidad de evaluar la carta} & E-ENT1-01. & la decisión comercial se vincula al costo y a la aceptación del precio. Por ello, conocer solo las ventas no basta para comparar rentabilidad por producto. & \textbf{Requisito:} R-ENT1-02, R-ENT1-03 y R-ENT1-04. \textbf{Conclusión:} necesidad de información para evaluar desempeño de la carta. \textbf{Límite:} el dato no define recetas, costos unitarios ni un algoritmo de margen por plato. \\
\textbf{I-ENT1-02. Tratamiento del antecedente} & E-ENT1-02. & la afirmación explica una decisión pasada. No obliga al sistema a prohibir productos artesanales ni demuestra que esa producción sea siempre inviable. & \textbf{Requisito:} ninguno. Conclusión de contexto: no se especifica una obligación independiente. \textbf{Conclusión:} antecedente del problema y pregunta sobre estructura de costo por producto. \textbf{Límite:} mantenerlo como contexto hasta obtener cantidades, rendimientos y costos verificables. \\
\textbf{I-ENT1-03. Corroboración del desencadenante} & E-ENT1-03. & la respuesta coincide con E-OBS-02. Esta convergencia documental refuerza el modelo de preparación por pedido, pero no acredita tiempos de preparación ni contempla todas las excepciones. & \textbf{Requisito:} ninguno propio. Corrobora R-OBS-02 (candidato de otro evento). \textbf{Conclusión:} regla candidata de iniciar la preparación tras recibir el pedido. \textbf{Límite:} confirmar tratamiento de bebidas listas, pedidos anticipados y productos especiales. \\
\textbf{I-ENT1-04. Registro de insumos y proveedores} & E-ENT1-04. & el negocio adquiere insumos con distintas presentaciones. El registro debe poder identificar lo comprado y su proveedor. Don Pollo es un valor del catálogo, no una constante del algoritmo. & \textbf{Requisito:} R-ENT1-01. \textbf{Conclusión:} necesidad de registrar proveedor, insumo y presentación de compra. \textbf{Límite:} la respuesta no establece exclusividad, frecuencia semanal ni integración digital con el proveedor. \\
\textbf{I-ENT1-05. Diferencia entre producto e insumo} & E-ENT1-05. & una bebida comprada lista y un refresco preparado requieren representación distinta si se desea calcular consumo o costo unitario. La descripción por sí sola no define esa representación. & \textbf{Requisito:} ninguno. Conclusión de contexto: no se especifica una obligación independiente. \textbf{Conclusión:} distinción conceptual entre producto de venta e insumo de preparación. \textbf{Límite:} no se infiere automáticamente un módulo de recetas ni control sanitario. \\
\end{longtable}

\subsection{Cadenas del evento EV-ENT-02}
\label{subsec:ilacion-ent2}

{\small
\begin{longtable}{L{2.35cm} L{2.15cm} L{4.35cm} L{5.15cm}}
\caption{Cadenas de ilación del evento EV-ENT-02 (28 cadenas).}
\label{tab:i-ent2}\\
\toprule
\textbf{Cadena} & \textbf{Evidencia} & \textbf{Razonamiento} & \textbf{Resultado y validación} \\
\midrule
\endfirsthead
\toprule
\textbf{Cadena} & \textbf{Evidencia} & \textbf{Razonamiento} & \textbf{Resultado y validación} \\
\midrule
\endhead
\bottomrule
\endfoot
\textbf{I-ENT2-01. Acceso dentro y fuera del negocio} & E-ENT2-01. & las ubicaciones expresan acceso remoto, no una arquitectura específica. & \textbf{Requisito:} R-ENT2-01. \textbf{Conclusión:} consulta administrativa remota autorizada. \textbf{Límite:} validar conectividad y funcionamiento ante pérdida de red. \\
\textbf{I-ENT2-02. Separación del trabajo administrativo} & E-ENT2-02. & distinguir tareas evita cargar la pantalla de caja con análisis gerencial. & \textbf{Requisito:} R-ENT2-01. \textbf{Conclusión:} vistas según tarea y entorno. \textbf{Límite:} separar vistas no exige dos aplicaciones o dos bases de datos. \\
\textbf{I-ENT2-03. Roles de uso} & E-ENT2-03. & los roles requieren permisos definidos, no acceso idéntico. & \textbf{Requisito:} R-ENT2-01. \textbf{Conclusión:} matriz de autorización por función. \textbf{Límite:} la respuesta no prohíbe al contador ver ventas; no mantener esa restricción sin acuerdo. \\
\textbf{I-ENT2-04. Interacción táctil y mouse} & E-ENT2-04. & cada contexto justifica interacción distinta. El tamaño de controles requiere prueba del dispositivo. & \textbf{Requisito:} R-ENT2-02. \textbf{Conclusión:} POS táctil y vista administrativa compatible con mouse. \textbf{Límite:} más de 48 px es un parámetro del equipo, no una medida indicada en la respuesta. \\
\textbf{I-ENT2-05. Resultado al final de la jornada} & E-ENT2-05. & la intención es obtener el resultado por una acción directa; no obliga a utilizar exclusivamente la tecla Enter. & \textbf{Requisito:} R-ENT2-03. \textbf{Conclusión:} resumen diario iniciado por el usuario. \textbf{Límite:} acordar horario de corte y tipos de venta incluidos. \\
\textbf{I-ENT2-06. Selección visual y comanda} & E-ENT2-06. & la comanda debe reflejar lo seleccionado. La analogía con McDonald's no exige copiar su interfaz. & \textbf{Requisito:} R-ENT2-02. \textbf{Conclusión:} pedido desglosado y comanda asociada. \textbf{Límite:} precisar impresión o monitor y momento de emisión. \\
\textbf{I-ENT2-07. Periodos y horarios ajustables} & E-ENT2-07. & los periodos de consulta y el corte operativo son conceptos diferentes. & \textbf{Requisito:} R-ENT2-03 y R-ENT2-04. \textbf{Conclusión:} filtros por periodo y jornada configurable. \textbf{Límite:} no se presume que la jornada coincide siempre con el día calendario. \\
\textbf{I-ENT2-08. Historial de precios} & E-ENT2-08. & aplicado a precios, permite conservar el valor cobrado aunque cambie la carta. & \textbf{Requisito:} ninguno propio. Corrobora R-PRE-13 (candidato de otro evento). \textbf{Conclusión:} precio histórico por venta y cambios identificados. \textbf{Límite:} definir permiso de cambio y vigencia. \\
\textbf{I-ENT2-09. Borrador y venta procesada} & E-ENT2-09. & una venta confirmada no debe borrarse como un borrador; su ajuste exige política específica. & \textbf{Requisito:} R-ENT2-14. \textbf{Conclusión:} eliminación de borrador y corrección trazable de venta confirmada. \textbf{Límite:} la frase nota de pedido no define todavía el documento de ajuste. \\
\textbf{I-ENT2-10. Búsqueda histórica} & E-ENT2-10. & la búsqueda exige conservar fecha e identidad estable del producto. & \textbf{Requisito:} R-ENT2-04. \textbf{Conclusión:} consulta por fecha y producto. \textbf{Límite:} distinguir productos homónimos de categorías diferentes. \\
\textbf{I-ENT2-11. Producto agotado} & E-ENT2-11. & la disponibilidad debe reflejarse en la selección; requiere un origen de datos de stock definido. & \textbf{Requisito:} R-ENT2-05. \textbf{Conclusión:} aviso o bloqueo del producto no disponible. \textbf{Límite:} no se deduce consumo automático por receta sin modelar insumos. \\
\textbf{I-ENT2-12. Modalidad de asignación elegida} & E-ENT2-12. & el tipo no se deriva solo del monto ni del proveedor. El usuario selecciona modalidad. & \textbf{Requisito:} R-ENT2-06. \textbf{Conclusión:} compra con modalidad diaria o semanal elegida. \textbf{Límite:} validar permisos y correcciones de modalidad. \\
\textbf{I-ENT2-13. Emergencia con matices} & E-ENT2-13. & el cargo directo es el caso habitual; la respuesta no acepta una imposición rígida para toda compra de emergencia. & \textbf{Requisito:} R-ENT2-06. \textbf{Conclusión:} valor por defecto diario con cambio autorizado según política. \textbf{Límite:} no inventar umbrales de monto que la respuesta no aporta. \\
\textbf{I-ENT2-14. Días hábiles y compras personales} & E-ENT2-14. & separar P-14a sobre calendario de P-14b sobre compras de empresa conserva ambas decisiones. & \textbf{Requisito:} R-ENT2-07. \textbf{Conclusión:} excluir jornadas no laboradas y separar gastos personales. \textbf{Límite:} la respuesta no solicita una función de porcentajes para facturas mixtas. \\
\textbf{I-ENT2-15. Cambio del precio de venta} & E-ENT2-15. & se necesita edición de precios, pero no se deriva ajuste automático ni garantía de no perder. & \textbf{Requisito:} ninguno propio. Corrobora R-PRE-12 y R-PRE-13. \textbf{Conclusión:} cambio manual autorizado con vigencia e historial. \textbf{Límite:} no recalcular precios de ventas anteriores. \\
\textbf{I-ENT2-16. Corrección de compra errónea} & E-ENT2-16. & la corrección debe ajustar asignaciones relacionadas y dejar rastro. & \textbf{Requisito:} R-ENT2-08 y R-ENT2-09. \textbf{Conclusión:} compra corregible con motivo y efecto controlado. \textbf{Límite:} P-24 limita modificaciones directas sobre periodos cerrados. \\
\textbf{I-ENT2-17. Plantilla de costos fijos} & E-ENT2-17. & se pueden reutilizar conceptos periódicos sin asumir que el monto permanece fijo. & \textbf{Requisito:} ninguno propio. Corrobora R-PRE-05. \textbf{Conclusión:} plantilla mensual editable de gastos indirectos. \textbf{Límite:} la financiación se valida aparte del esquema verbal. \\
\textbf{I-ENT2-18. Norma y autorización del dueño} & E-ENT2-18. & distinguir referencia normativa de autorización operativa. Una opinión sobre SUNAT requiere contraste externo. & \textbf{Requisito:} R-ENT2-10. \textbf{Conclusión:} cambios de configuración contable con autorización del dueño. \textbf{Límite:} la declaración no acredita obligatoriedad legal de NIIF 18 para este negocio. \\
\textbf{I-ENT2-19. Vigencia de cambios tributarios} & E-ENT2-19. & la regla aplicable se identifica por fecha de efecto, no por la fecha de edición. & \textbf{Requisito:} R-ENT2-10 y R-ENT2-11. \textbf{Conclusión:} versiones tributarias con inicio de vigencia. Fuente puntual: N-05, Título Preliminar, Normas IV y X; consultar la disposición de vigencia de cada modificación. \textbf{Límite:} no afirmar que ningún cambio a mitad de mes pueda ocurrir; verificar norma concreta. \\
\textbf{I-ENT2-20. Aprobación del reporte} & E-ENT2-20. & la plantilla de presentación debe contar con revisión del responsable contable. & \textbf{Requisito:} ninguno propio. Corrobora R-PRE-08. \textbf{Conclusión:} reporte mensual con versión aprobada. \textbf{Límite:} no tratar el modelo inicial como plantilla final certificada. \\
\textbf{I-ENT2-21. Prioridad de contenido del reporte} & E-ENT2-21. & el sistema conserva detalle para auditoría, pero la vista principal debe entregar el estado mensual. & \textbf{Requisito:} ninguno propio. Corrobora R-PRE-08. \textbf{Conclusión:} vista mensual con acceso al sustento cuando se solicite. \textbf{Límite:} no imponer una pantalla tutorial de cálculo como requisito del cliente. \\
\textbf{I-ENT2-22. Archivo electrónico ligado al gasto} & E-ENT2-22. & el problema es localizar el sustento, no digitalizar obligatoriamente un papel. & \textbf{Requisito:} R-ENT2-12. \textbf{Conclusión:} vínculo egreso-comprobante y consulta desde reportes. \textbf{Límite:} permitir cargar archivo existente; foto de celular no es modalidad obligatoria. \\
\textbf{I-ENT2-23. Referencia a Camaleón} & E-ENT2-23. & la mención sirve para preguntar por expectativas de integración; no demuestra homologación del sistema propuesto. & \textbf{Requisito:} ninguno. Conclusión de contexto: no se especifica una obligación independiente. \textbf{Conclusión:} cuestión de alcance sobre respaldos e integración. \textbf{Límite:} verificar producto y obligación concreta antes de especificar homologación. \\
\textbf{I-ENT2-24. Ajuste de periodo cerrado} & E-ENT2-24. & puede conservarse el cierre y permitir una corrección por procedimiento controlado. & \textbf{Requisito:} R-ENT2-09, R-ENT2-13 y R-ENT2-14. \textbf{Conclusión:} periodo cerrado sin edición directa y ajuste autorizado trazable. \textbf{Límite:} definir autorización, tipo de nota y periodo del efecto. \\
\textbf{I-ENT2-25. Compresión de adjuntos} & E-ENT2-25. & la reducción no debe perder legibilidad ni desvincular el archivo. & \textbf{Requisito:} R-ENT2-15. \textbf{Conclusión:} adjunto optimizado y recuperable. \textbf{Límite:} aprobar límite de tamaño, formatos y preservación del original. \\
\textbf{I-ENT2-26. Validación de entradas} & E-ENT2-26. & se requieren validaciones específicas por campo; algunos ajustes pueden necesitar importes negativos autorizados. & \textbf{Requisito:} R-ENT2-16. \textbf{Conclusión:} mensaje explicativo y rechazo de entrada inválida. \textbf{Límite:} definir reglas por operación y excepciones de ajustes. \\
\textbf{I-ENT2-27. Emisión SUNAT y expectativa de afiliación} & E-ENT2-27. & la respuesta confirma emisión externa con una expectativa no resuelta. Debe consignarse como acuerdo condicionado. & \textbf{Requisito:} ninguno. Conclusión de contexto: no se especifica una obligación independiente. \textbf{Conclusión:} emisión fuera del sistema; integración pendiente de decisión documentada. \textbf{Límite:} la opinión no constituye una obligación legal verificada. \\
\textbf{I-ENT2-28. Ventas por distintos canales} & E-ENT2-28. & registrar el canal es distinto de controlar repartidores. & \textbf{Requisito:} R-ENT2-17. \textbf{Conclusión:} ventas de distintos canales en el mismo registro. \textbf{Límite:} no inferir GPS, ruteo o administración de motorizados. \\
\end{longtable}
}

\subsection{Integración y trazabilidad entre eventos}
\label{subsec:integracion}

La integración conecta conclusiones ya identificadas. La secuencia temporal
de operación y abastecimiento no implica que una entrevista haya producido
todos los requisitos del sistema. La Tabla~\ref{tab:trazabilidad} conserva
tanto la trazabilidad hacia el origen como la relación hacia el
comportamiento especificado.

\begin{longtable}{L{4.2cm} L{4.5cm} L{5.6cm}}
\caption{Integración de resultados y trazabilidad entre eventos.}
\label{tab:trazabilidad}\\
\toprule
\textbf{Resultado integrado} & \textbf{Datos de origen} & \textbf{Ilación y especificación} \\
\midrule
\endfirsthead
\toprule
\textbf{Resultado integrado} & \textbf{Datos de origen} & \textbf{Ilación y especificación} \\
\midrule
\endhead
\bottomrule
\endfoot
Comanda y preparación por pedido & E-OBS-02; E-ENT1-03; D-COM-02 & I-OBS-02 y análisis de D-COM-02 deriva en R-OBS-02. \\
Compra y proveedor & E-ENT1-04 & I-ENT1-04 deriva en R-ENT1-01. \\
Clasificación y reparto del costo & E-PRE-05; contraste con F-02, 5.2 & I-PRE-05 deriva en R-PRE-01/02/03. \\
Margen diario y resultado mensual & E-PRE-04/05/06/07/14 & I-PRE-04/05/06/07 deriva en R-PRE-04/05/06. \\
Comprobante asociado & E-PRE-11/13 & I-PRE-11/13 deriva en R-PRE-09. \\
Reporte contable y consulta remota & E-PRE-10/14; D-COM-01/03 & I-PRE-10/14 deriva en R-PRE-08; I-ENT2-01/02 deriva en R-ENT2-01. \\
\end{longtable}

\paragraph{Contradicciones que requieren validación.} El reporte de origen
identifica seis diferencias que no deben resolverse por invención:
(i) cinco días de la propuesta inicial frente a días hábiles efectivos de
F-08 P-14; (ii) el BAI que resta financiación después de un subtotal así
denominado; (iii) los permisos del contador (F-02 los limita, F-08 P-03 no
expresa esa prohibición); (iv) SHA-256 presentado como cifrado cuando es
hash de integridad; (v) el rollback como respuesta completa ante corte
eléctrico; y (vi) la equiparación entre reporte gerencial, NIIF 18 y SUNAT.
Todas permanecen abiertas y se trasladan a la Sección~\ref{sec:especificacion}.
